\scanpage{71}
\uline{Variante} laissons tomber l'hypothèse $d'\ (=\mathrm{pgcd}\text{ des deg. des 0-cycles sur }X_\eta)=1$.
On sait que $\operatorname{Br}(K)$ est un gpe de $p$-torsion, donc \unsure{nul|on a}.
\[
\unsure{0|\cdot}\longrightarrow\operatorname{Pic}(X_\eta)\mathbin{-}\unsure{P|R}(K)\mathbin{-}\Theta\mathbin{-}0,
\]
où $\Theta\subset\operatorname{Br}(K)$ est un gpe de $p$-torsion, d'où résulte que ${}_n\operatorname{Pic}(X_\eta)\simeq{}_nP(K)$. Donc il reste vrai que $u_n$ est isom. à $\unsure{{}_nv|n^0|u^0}$. Mais $\Gamma$ peut avoir de la torsion. On trouve
\diagram{\begin{tikzcd}[column sep=small]
0\arrow[r]&{}_n\Gamma\arrow[r]\arrow[d,phantom,"\unsure{\simeq}"{description,rotate=90}]&\mathrm H^1(X,\mu_n)\arrow[r,dash]&\mathrm H^1(X_{\bar\eta},\mu_n)^I\arrow[r]&\Gamma_n\arrow[r,dash]&\operatorname{NS}(X_0)_n\arrow[r]&{}\\
&\mathbf Z/(n,d)\mathbf Z
\end{tikzcd}}
\deleted{Soit} et par passage à la limite sur les puissances de $\ell$
\diagram{\begin{tikzcd}[column sep=small]
0\arrow[r,dash]&\mathrm H^1(X,\mathbf Z_\ell(1))\arrow[r]&\mathrm H^1(X_{\bar\eta},\mathbf Z_\ell(1))^I\arrow[r]\arrow[dr,rounded corners,to path={|- (\tikztotarget) \tikztonodes}]&\Gamma\otimes\mathbf Z_\ell\arrow[r,dash]&\operatorname{NS}(X_0)\otimes\mathbf Z_\ell\\
&&&R\otimes\mathbf Z_\ell\arrow[r,dash]&0,
\end{tikzcd}}
$R=\operatorname{Ker}(\Gamma\mathbin{-}\operatorname{NS}(X_0))$
\margin{N.B. $R\otimes\mathbf Z_\ell\simeq\begin{cases}\mathbf Z/\ell\mathbf Z&\text{si }\ell\mid d\\0&\text{sinon}\end{cases}$}
On voit ainsi (tout au moins modulo rés. des singularités, pour pouvoir faire \unsure{qch.|cas.} de « bon) que le noyau $R$ de $\Gamma\mathbin{-}\operatorname{NS}(X_0)$ est de \uline{torsion}. \deleted{(\unsure{donc} un gpe cyclique \illegible{1} d'ordre divisant $d$, $d\,\unsure{?}\,d'$)} C'est donc un sous gpe de $\operatorname{tors}\Gamma\simeq\mathbf Z/d\mathbf Z$, [avec une nette tendance d'être égal à $\mathbf Z/d\mathbf Z$ tout entier ; \deleted{\unsure{et}} tout au moins le noyau contient-il la partie de $\operatorname{tors}\Gamma$ première à $p$, pt argument de Raynaud utilise un rev. étale convenable de $S$ ....) On a donc, si
\margin{\unsure{même|vrai} sans résolution \unsure{référer|référence} à Raynaud, note citée, prop. 7 (ii)}
\margin{demander Raynaud si $\operatorname{Ker}$ est \unsure{donc|tout|bien} $\mathbf Z/d\mathbf Z$.}
\[
\Gamma'=\Gamma/T\deleted{\unsure{orsion}},\quad\deleted{\unsure{\&|R}}\quad T=\operatorname{Torsion}\text{ de }\Gamma
\]
\[
\unsure{0|\cdot}\mathbin{-}T_n\mathbin{-}\Gamma_n\mathbin{-}\Gamma'_n\longrightarrow0
\]
d'où, utilisons un splitting de $\Gamma\mathbin{-}\Gamma'\oplus T$, on trouve
\[
\operatorname{Ker}(\Gamma_n\mathbin{-}\operatorname{NS}(X_0)_n)\simeq T_n+\underbrace{\operatorname{Ker}(\Gamma'_n\mathbin{-}\operatorname{NS}(X_0)_n)}_{\operatorname{Ker}(\overline\Gamma_n\mathbin{-}\operatorname{NS}(X_0)_n)}
\]
où
\diagram{\begin{tikzcd}[column sep=small,row sep=small]
\overline\Gamma\arrow[r,hook]\arrow[d,equal]&\operatorname{NS}(X_0)\\
\operatorname{Im}\Gamma
\end{tikzcd}}
On trouve donc, comme ${}_n\overline\Gamma$ \deleted{\unsure{est}} n'a que de la $p$-torsion
\diagram{\begin{tikzcd}[column sep=large]
0\arrow[r,dash]&\delta'\arrow[r]\arrow[d,equal]&\delta\arrow[r,dash]&\delta''\arrow[r,dash]&0\\
&\text{tors de }\operatorname{NS}(X_0)\text{ \unsure{première|premier} à }p
\end{tikzcd}}
\margin{70}
